\documentclass{article}
\usepackage{amsmath}
\usepackage{graphicx} % Required for inserting images


\title{Modello Big project}
\begin{document}


\section{Set}
\begin{itemize}
    \item $N$: Insieme di tutti i nodi (punti griglia $S$ e candidati base $B$).
    \item $S$: Insieme dei nodi da visitare.
    \item $B$: Insieme dei possibili punti di base ($B \subset N$).
    \item $H$: Insieme dei viaggi disponibili, indicizzati da $h$. 
    \\$|H|$  va scelta correttamente: 
    il numero massimo di viaggi deve essere sufficiente a garantire la copertura di tutti i punti, ma limitato per contenere l'esplosione delle variabili decisionali $x_{ij}^h$. 
        \begin{equation}
|H| = \left\lceil \frac{\sum_{i \in S} (e_{in}^{min} + e_{out}^{min})}{C} \cdot 1.5 \right\rceil
\end{equation}
    dove $e_{in}^{min}$ e $e_{out}^{min}$ sono gli archi con il costo minimo per entrare e uscire dal punto $i$.
    \item $A$: Insieme degli archi validi $(i,j)$ che rispettano i vincoli di distanza e connessione.
\end{itemize}
\section{Parametri}
\begin{itemize}
    \item $t_{ij}$: Tempo di percorrenza sull'arco $(i,j)$.
    \item $e_{ij}$: Consumo energetico sull'arco $(i,j)$.
    \item $C$: Capacità massima della batteria.
    \item $M$: Costante grande necessaria per i subtour.
\end{itemize}
\section{Variabili}
\begin{align*}
    x_{ij}^h & \in \{0, 1\} && \text{1 se il viaggio } h \text{ percorre l'arco } (i,j), \text{ 0 altrimenti.}\\
    y_b & \in \{0, 1\} && \text{1 se il nodo } b \in B \text{ è scelto come base, 0 altrimenti.} \\
    f_{ij}^h & \ge 0 && \text{Variabile di flusso per l'eliminazione dei subtour (SCF).}
\end{align*}

\section{Formulazione}

\subsection{Funzione Obiettivo:}
\begin{equation}
    \min \sum_{h \in H} \sum_{(i,j) \in A} t_{ij} \cdot x_{ij}^h
\end{equation}
\textit   Minimizzare il tempo totale di volo necessario per completare la copertura.

\subsection{Vincoli}

\paragraph{1. Selezione Unica della Base}
\begin{equation}
    \sum_{b \in B} y_b = 1
\end{equation}
\textit   Deve essere installata una sola stazione di ricarica in un unico punto $(x_0, y_0, z_0)$ valido per tutti i viaggi.

\paragraph{2. Copertura dei Punti}
\begin{equation}
    \sum_{h \in H} \sum_{j \in N} x_{ji}^h \ge 1 \quad \forall i \in S
\end{equation}
\textit   Ogni punto della griglia deve essere visitato almeno una volta complessivamente sui viaggi h.

\paragraph{3. Conservazione del Flusso}
\begin{equation}
    \sum_{j \in N} x_{ji}^h - \sum_{j \in N} x_{ij}^h = 0 \quad \forall i \in N, \forall h \in H
\end{equation}
\textit   Se il drone entra in un nodo durante un viaggio, deve necessariamente uscirne.

\paragraph{4. Coerenza Base-Viaggio }
\begin{equation}
    \sum_{j \in S} x_{bj}^h \le y_b \quad \forall b \in B, \forall h \in H
\end{equation}
\textit    Se $y_b=0$ (nodo non scelto come base), nessun viaggio può partire da $b$. Se $y_b=1$, il viaggio $h$ può partire da $b$ al massimo una volta.
 L'uso del minore o uguale garantisce l'ammissibilità del problema poiché permette al solvitore di attivare solo i viaggi strettamente necessari evitando "viaggi vuoti" forzati.
  Non vi è rischio di percorsi isolati fuori dalla base poiché la conservazione del flusso e il vincolo MTZ impediscono matematicamente la formazione di cicli chiusi che non includano il punto di partenza b .

\paragraph{5. Capacità della Batteria}
\begin{equation}
    \sum_{(i,j) \in A} e_{ij} \cdot x_{ij}^h \le C \quad \forall h \in H
\end{equation}
\textit   L'energia totale consumata in ogni singolo viaggio $h$ non può superare la capacità della batteria $C$. Il drone si ricarica completamente tra un viaggio e l'altro.

\paragraph{6. Eliminazione dei Subtour (SCF)}

Per prevenire la formazione di subtour isolati nei percorsi dei veicoli, utilizziamo un approccio basato sul \textbf{Single Commodity Flow (SCF)}.  


L'idea principale è che un flusso parte dalla base e viene consumato ad ogni nodo visitato: se esiste un sottotour che non passa per la base, non avrà sorgente di flusso e diventa unfeasible.




\begin{enumerate}
    \item \textbf{Flusso attivo solo sugli archi selezionati:}
    \begin{equation}
        f_{ij}^h \le |S| \cdot x_{ij}^h \quad \forall (i,j) \in A, \forall h \in H
    \end{equation}
    Questo vincolo garantisce che il flusso possa circolare solo sugli archi attivi, cioè quelli effettivamente selezionati nel percorso.

    \item \textbf{Conservazione del flusso nei nodi di servizio:}
    \begin{equation}
        \sum_{j \in N} f_{ji}^h - \sum_{j \in N} f_{ij}^h = \sum_{j \in N} x_{ji}^h \quad \forall i \in S, \forall h \in H
    \end{equation}
    La differenza tra flusso in ingresso e in uscita da un nodo di servizio deve corrispondere al numero di volte in cui il nodo viene visitato. Questo impedisce che i nodi vengano visitati senza che il flusso arrivi dalla base.

    \item \textbf{Flusso alle basi:}  
    La base funge da sorgente primaria. Non è necessario introdurre un vincolo esplicito perché i vincoli di conservazione del flusso (2) assicurano che il flusso totale consumato sia pari al numero di nodi visitati, implicando che ogni percorso inizi dalla base.
\end{enumerate}

\paragraph{Osservazioni:}

\begin{itemize}
    \item Ogni tour deve partire e tornare alla base: un ciclo chiuso tra soli nodi di servizio è impedito dal vincolo di flusso, così come un cammino aperto o un tour che non ritorna alla base.
    \item L'approccio SCF riduce il numero di vincoli rispetto all'esplicita enumerazione dei subtour.
\end{itemize}


\paragraph{7. Ordinamento dei Viaggi (Symmetry Breaking)}
\begin{equation}
    \sum_{(i,j) \in A} x_{ij}^{h} \ge \sum_{(i,j) \in A} x_{ij}^{h+1} \quad \forall h \in \{1, \dots, |H|-1\}
\end{equation}
\textit   Forza il solvitore a utilizzare i viaggi in ordine crescente in modo da ridurre lo spazio di ricerca. Senza questo constraint, ci sono soluzioni simmetriche.
\end{document}
